% Pandoc template for IEEE conference paper sections
\section{Curve1024: Parameters and Evaluation}

This section presents the concrete Curve1024 parameters produced by our
improved Cocks-Pinch construction with readability-aware selection,
followed by implementation details and security analysis.

\subsection{Curve1024 Parameters}\label{curve1024-parameters}

The construction produces the following pairing-friendly elliptic curve
over a 1024-bit prime field. All parameters are publicly verifiable from
the seed \(T\).

\subsubsection{Seed and Construction}\label{seed-and-construction}

The curve is fully determined by a single 86-bit seed:

\[T = \texttt{0x26704d2ace0facdd539} \cdot 2^{12}\]

This seed satisfies \(T \equiv 0 \pmod{2^{12}}\), guaranteeing
\(\nu_2(r-1) = 36\) via the factorization \(r - 1 = T^3(T^3 - 1)\).

\subsubsection{Scalar Field Order (512
bits)}\label{scalar-field-order-512-bits}

The scalar field prime \(r = \Phi_{18}(T) = T^6 - T^3 + 1\) decomposes
as:

\[r = d_r \cdot 2^{36} + 1\]

where \(d_r\) is a 476-bit odd integer. In 64-bit limb representation
(MSB first):

\begin{table}[t]
\centering
\caption{Scalar field prime \(r\) --- limb decomposition}
\begin{tabular}{lll}
\toprule
Limb & Value & Note \\
\midrule
15 & \texttt{0000000000000000} & zero \\
14 & \texttt{0000000000000000} & zero \\
13 & \texttt{0000000000000000} & zero \\
12 & \texttt{0000000000000000} & zero \\
11 & \texttt{0000000000000000} & zero \\
10 & \texttt{0000000000000000} & zero \\
9 & \texttt{0000000000000000} & zero \\
8 & \texttt{0000000000000000} & zero \\
7 & \texttt{c042c3389e72044d} &  \\
6 & \texttt{9ec8078aea7bc954} &  \\
5 & \texttt{ebf209d918f0ef27} &  \\
4 & \texttt{425259e247711927} &  \\
3 & \texttt{97463bf23832d0de} &  \\
2 & \texttt{5e46e47fc61cab50} &  \\
1 & \texttt{4bcfd36fc2a945c5} &  \\
0 & \texttt{1055d97000000001} & tail: \(d \cdot 2^{36} + 1\) \\
\bottomrule
\end{tabular}
\end{table}

\textbf{Key readability features:} 8 of 16 limbs are zero (50\%), and
the final limb visually confirms the NTT-friendly structure with the
\texttt{00000001} suffix (36 trailing zero bits before the \(+1\)).

\subsubsection{Base Field Prime (1024
bits)}\label{base-field-prime-1024-bits}

The base field prime \(p = (t^2 + 3y^2)/4\) decomposes as:

\[p = d_p \cdot 2^{36} + 1\]

where \(d_p\) is a 988-bit odd integer. The final limb shares the same
tail pattern as \(r\):

\begin{table}[t]
\centering
\caption{Base field prime \(p\) --- limb decomposition}
\begin{tabular}{lll}
\toprule
Limb & Value & Note \\
\midrule
15 & \texttt{c0859da83a500988} &  \\
14 & \texttt{10eea35c61fea054} &  \\
13 & \texttt{c866902599220d30} &  \\
12 & \texttt{c93f38265ed6b830} &  \\
11 & \texttt{5785b8a073e482e8} &  \\
10 & \texttt{685cfeeb388dc45e} &  \\
9 & \texttt{bca1686245b326fe} &  \\
8 & \texttt{824ca74ac0844160} &  \\
7 & \texttt{f4e289806084e50e} &  \\
6 & \texttt{6c02f32cb94d786b} &  \\
5 & \texttt{f80fb080c07d6a07} &  \\
4 & \texttt{73f2b4591eb70df3} &  \\
3 & \texttt{9367aa32ebb42449} &  \\
2 & \texttt{e0a648f61a20e876} &  \\
1 & \texttt{f971f3bf976a7d6a} &  \\
0 & \texttt{dd20d97000000001} & tail: \(d \cdot 2^{36} + 1\) \\
\bottomrule
\end{tabular}
\end{table}

\textbf{Shared tail pattern:} Both \(p\) and \(r\) end with
\texttt{...d970\_00000001}, a direct consequence of the seed constraint
propagating through the construction. This visual consistency aids
manual verification.

\subsubsection{Curve Summary}\label{curve-summary}

\begin{table}[t]
\centering
\caption{Curve1024 parameters}
\begin{tabular}{ll}
\toprule
Parameter & Value \\
\midrule
Equation & \(E: y^2 = x^3 + 41\) \\
Base field \(|p|\) & 1024 bits \\
Scalar field \(|r|\) & 512 bits \\
Embedding degree \(k\) & 18 \\
CM discriminant \(D\) & \(-3\) (hence \(j = 0\)) \\
\(\rho = \log p / \log r\) & 2.0 \\
\(\nu_2(p-1)\) & 36 \\
\(\nu_2(r-1)\) & 36 \\
Cofactor \(h = \#E/r\) & \(\approx 2^{512}\) \\
\bottomrule
\end{tabular}
\end{table}

The curve equation \(y^2 = x^3 + 41\) arises from the CM method with
\(D = -3\): the Hilbert class polynomial \(H_{-3}(x) = x\) gives
\(j = 0\), yielding the family \(y^2 = x^3 + b\). The twist \(b = 41\)
is selected by verifying which of the six sextic twists has group order
divisible by \(r\).

A generator point \(G = (G_x, G_y)\) of prime order \(r\) is obtained by
cofactor multiplication of a random curve point. The full coordinates
are published in the project repository \cite{LumenMath}.

\subsubsection{Readability Metrics}\label{readability-metrics}

\begin{table}[t]
\centering
\caption{Readability comparison}
\begin{tabular}{lll}
\toprule
Metric & \(p\) (1024-bit) & \(r\) (512-bit) \\
\midrule
Hamming weight & 464/1024 (45\%) & 235/512 (46\%) \\
Two-adicity & 36 & 36 \\
Zero 64-bit limbs & 0/16 & 8/16 (50\%) \\
Hex zero digits & 33/256 (13\%) & 16/128 (12.5\%) \\
Trailing zero bits & 36 & 36 \\
\bottomrule
\end{tabular}
\end{table}

The scalar field prime \(r\) is notably sparse: half its limb
representation is zero, making it easy to recognize and verify in
debugging output. Both primes exceed BLS12-381's two-adicity of 32.

\subsection{Implementation}\label{implementation}

We implement the complete system in Rust without external cryptographic
dependencies \cite{LumenMath}. The implementation comprises:

\begin{itemize}
\tightlist
\item
  A custom \texttt{U1024} big integer type with Montgomery
  multiplication \cite{Montgomery1985}
\item
  Prime field arithmetic (\texttt{PrimeFieldElement}) with constant-time
  Fermat inversion
\item
  Elliptic curve group operations in affine coordinates
\item
  Schnorr and ECDSA signature schemes
\item
  The parameter generation tool (improved Cocks-Pinch + readability
  scoring)
\end{itemize}

The test suite contains 36 tests covering arithmetic correctness, group
law verification, and attack simulation (MOV, Anomalous, Invalid Curve).

\subsection{Performance Benchmarks}\label{performance-benchmarks}

\begin{table}[t]
\centering
\caption{Performance benchmarks (Criterion, release build, LLVM O3, no AVX-512)}
\begin{tabular}{llll}
\toprule
Operation & Min & Mean & Max \\
\midrule
Key generation & 1.629 s & 1.687 s & 1.757 s \\
Schnorr sign & 1.473 s & 1.551 s & 1.641 s \\
Schnorr verify & 2.656 s & 2.831 s & 3.025 s \\
ECDSA sign & 1.524 s & 1.606 s & 1.711 s \\
ECDSA verify & 2.847 s & 3.140 s & 3.459 s \\
Scalar mult. \([k]G\) & 1.061 s & \textasciitilde1.2 s & 1.489 s \\
\bottomrule
\end{tabular}
\end{table}

The performance reflects the 1024-bit field size: each scalar
multiplication requires \textasciitilde512 double-and-add iterations,
each involving a Fermat inversion (\textasciitilde1536 Montgomery
multiplications). This is approximately \(8000\times\) slower than
optimized secp256k1 implementations, which is expected given the
\(4\times\) larger field and \(O(n^2)\) multiplication cost. The system
targets offline signing scenarios (legal documents, software releases,
PKI certificates) where seconds of latency are acceptable.

\subsection{Security Analysis}\label{security-analysis}

\begin{table}[t]
\centering
\caption{Security evaluation summary}
\begin{tabular}{lll}
\toprule
Attack & Complexity & Status \\
\midrule
Pollard-\(\rho\) (ECDLP) & \(O(2^{256})\) & Secure \\
MOV reduction & DLP in \(\mathbb{F}_{p^{18}}\) (18432-bit) & Secure \\
Anomalous (SSSA) & Requires \(\#E = p\) & Immune \\
TNFS {[}@BarbulescuDuquesne2019{]} & \(\geq 250\)-bit in \(\mathbb{F}_{p^{18}}\) & Secure \\
Invalid curve & Point validation in constructor & Mitigated \\
\bottomrule
\end{tabular}
\end{table}

The 256-bit ECDLP security exceeds NIST SP 800-57 recommendations
through 2030+ \cite{NIST80057}. The MOV attack \cite{MOV1993} maps
ECDLP to DLP in \(\mathbb{F}_{p^{18}}\), a field of 18,432 bits---far
exceeding the 3,072-bit threshold for 128-bit security. The curve is
provably non-anomalous \cite{SmartAnomaly1999} since
\(\#E(\mathbb{F}_p) = h \cdot r\) with
\(r \approx 2^{512} \neq p \approx 2^{1024}\).

\textbf{Limitation:} The current affine-coordinate implementation uses
variable-time double-and-add, which leaks scalar bits through timing. A
Montgomery ladder (constant-time) is identified as future work.

\subsection{Curve Family Comparison}\label{curve-family-comparison}

\begin{table}[t]
\centering
\caption{Comparison with standard pairing-friendly curve families}
\begin{tabular}{llllll}
\toprule
Curve & \(|p|\) (bits) & \(\rho\) & \(\nu_2\) & ECDLP Security & Pairing Target \\
\midrule
BN254 & 254 & 1.0 & \textasciitilde1 & 100-bit & 3048-bit \\
BLS12-381 & 381 & 1.5 & 32 & 128-bit & 4572-bit \\
KSS18 (generic) & varies & 1.33 & random & varies & varies \\
\textbf{Curve1024} & \textbf{1024} & 2.0 & \textbf{36} & \textbf{256-bit} & \textbf{18432-bit} \\
\bottomrule
\end{tabular}
\end{table}

Curve1024 occupies a unique position: it provides the highest ECDLP
security level (256-bit) and the highest two-adicity (36) among deployed
pairing-friendly curves, at the cost of a larger \(\rho\)-value inherent
to the Cocks-Pinch method. The 18,432-bit pairing target field provides
substantial margin against TNFS attacks.
